\chapter{CONCLUSION AND RECOMMENDATION}\label{chap:conclusion-recommendation}

This chapter presents the conclusions and recommendations drawn from the
development and evaluation of the TinyML Smart Plug. The conclusions
summarize how the prototype addressed the study objectives through calibrated
sensing, threshold-based overload and socket-heating protection, TinyML-supported
series arc-fault detection, relay actuation, local notification, and PWA monitoring.
The recommendations identify the remaining validation, dataset, thermal,
hardware-safety, logging, and embedded-platform improvements needed before the
prototype can be considered for residential application.

\section{Conclusions}\label{sec:conclusions}

The study demonstrated that the TinyML Smart Plug can operate as a
complementary socket-level protection prototype under controlled laboratory
conditions. The calibrated voltage, current, and temperature channels produced
errors low enough to support the implemented protection decisions, with the
sensor-validation results reporting mean absolute errors of 3.910V, 0.178A, and
1.02°C, respectively. These results support the first objective because the
prototype was able to acquire the required electrical and thermal quantities and
use them in embedded decision logic for monitoring, warning, and relay control.

The overload tests showed that the prototype could distinguish warning-level
current from sustained overload. All overload trials generated warnings, while
only the trial with a 17.30A mean current produced a relay trip after 61.3s.
This agrees with the implemented persistence method and supports the objective
of adding time-aware overload response rather than using an instantaneous
cutoff for every current excursion. This is consistent with the reviewed
literature describing overload as a time-dependent current and thermal-stress
condition, where the protection response must distinguish short transients from
sustained heating risk [12]-[14], [22], [44], [45].

The socket-heating tests confirmed that degraded contact can create faster and
more severe heating than a normal socket at comparable current levels. The
estimated socket temperature reached the 40°C trip threshold earlier in
corroded-socket trials, and the reference terminal temperature was already
higher than the device temperature reading at trip. This supports the need for a
thermal protection layer that is separate from current-only protection. This
follows the reviewed literature showing that poor plug-socket contact can
produce localized and hidden heating even when the total current does not
necessarily force conventional upstream protection to operate [17], [24],
[46], [47]. It also supports the reviewed finding that external temperature
measurement can under-read the hottest contact region because of sensor
placement, thermal lag, and thermal separation from the internal hotspot
[27]-[31].

The series arc-fault tests showed that embedded TinyML inference was practical
on the implemented platform. Across the tested devices, the prototype achieved
97.50\% accuracy, 98.97\% precision, 96.00\% recall, and a 97.46\% F1-score.
These results support the objective of detecting arc-related waveform
disturbance using a feature-based embedded model. This is consistent with the
reviewed literature showing that series arcs may remain below ordinary breaker
trip levels while still producing conduction intermittence, waveform
disturbance, and spectral changes that require multi-feature interpretation
rather than RMS current magnitude alone [16], [33]-[39], [49]-[54]. The use of
load context also supports the reviewed finding that healthy nonlinear loads can
produce distortion that overlaps with arc-like behavior, making context-aware
inference important for reducing nuisance operation [8], [40]-[43].

The prototype is therefore best interpreted as a complementary protection
device for selected progressive fire-inducing fault conditions. It was able to
combine calibrated voltage, current, and temperature sensing with deterministic
overload and socket-heating protection, TinyML-supported series arc-fault
interpretation, relay actuation, notification, event logging, and PWA monitoring
in one working prototype. However, the results must be interpreted within the
tested load set, induced-fault procedures, and controlled laboratory conditions
only. The prototype is not a replacement for circuit breakers, AFCIs, GFCIs,
proper wiring practice, certified receptacles, inspection, or other mandatory
protection devices.

\section{Recommendations}\label{sec:recommendations}

Longer validation should be performed in residentially realistic or semi-field
conditions. Additional tests should include ordinary daily load transitions,
shared branch circuits, extension cords, aged receptacles, different plug-fit
conditions, and different wiring lengths.

The arc-fault dataset should be expanded to include more load families,
additional appliance models, more arc severities, and non-fault disturbances
such as motor startup, switching noise, loose plugs without confirmed arcing,
and intermittent contact. This would give the Random Forest model stronger
coverage of nuisance events and would help improve recall for less obvious
arc-fault cases.

The socket-temperature subsystem should be improved through better thermal
sensor placement and additional model validation. Subsequent prototypes should
evaluate thermistors placed closer to the plug-contact interface, multiple
temperature sensors, and longer degraded-contact trials so that the relationship
among device temperature measurement, estimated socket temperature, and actual
terminal temperature can be fitted with stronger evidence.

The enclosure, creepage and clearance distances, insulation barriers, relay
rating, wire gauge, strain relief, and thermal isolation should be redesigned
for certification-oriented testing. Additional tests should be aligned with
applicable electrical safety and appliance standards before any field or
consumer deployment is considered.

The monitoring system should add stronger event logging and data export.
Exportable CSV logs, synchronized timestamps across waveform and telemetry
events, firmware-version fields, and clearer fault-state records would make
later validation easier and would support model retraining from collected test
data.

Finally, later implementations may use a microcontroller platform with more
memory and data-handling margin. Additional resources would allow longer
waveform buffers, larger validation logs, broader feature comparison, and more
robust embedded inference while keeping the protection logic responsive.